%!TeX root=../treasuretop.tex
\ArtChapter[The Voyage]{chap10head}
 
   \lettrine[lines=4,image=true]{chap10initA}{ll} that night we were in a great bustle getting things stowed in their place, and boatfuls of the squire's friends, Mr Blandly and the like, coming off to wish him a good voyage and a safe return. We never had a night at the »Admiral Benbow« when I had half the work; and I was dog-tired when, a little before dawn, the boatswain sounded his pipe and the crew began to man the capstan-bars. I might have been twice as weary, yet I would not have left the deck, all was so new and interesting to me—the brief commands, the shrill note of the whistle, the men bustling to their places in the glimmer of the ship's lanterns.

»Now, Barbecue, tip us a stave,« cried one voice.

»The old one,« cried another.

»Aye, aye, mates,« said Long John, who was standing by, with his crutch under his arm, and at once broke out in the air and words I knew so well:

»Fifteen men on the dead man's chest\longdash«

And then the whole crew bore chorus:—

»Yo-ho-ho, and a bottle of rum!«

And at the third »Ho!« drove the bars before them with a will.

Even at that exciting moment it carried me back to the old »Admiral Benbow« in a second, and I seemed to hear the voice of the captain piping in the chorus. But soon the anchor was short up; soon it was hanging dripping at the bows; soon the sails began to draw, and the land and shipping to flit by on either side; and before I could lie down to snatch an hour of slumber the \textit{Hispaniola} had begun her voyage to the Isle of Treasure.

I am not going to relate that voyage in detail. It was fairly prosperous. The ship proved to be a good ship, the crew were capable seamen, and the captain thoroughly understood his business. But before we came the length of Treasure Island, two or three things had happened which require to be known.

Mr Arrow, first of all, turned out even worse than the captain had feared. He had no command among the men, and people did what they pleased with him. But that was by no means the worst of it, for after a day or two at sea he began to appear on deck with hazy eye, red cheeks, stuttering tongue, and other marks of drunkenness. Time after time he was ordered below in disgrace. Sometimes he fell and cut himself; sometimes he lay all day long in his little bunk at one side of the companion; sometimes for a day or two he would be almost sober and attend to his work at least passably.

In the meantime, we could never make out where he got the drink. That was the ship's mystery. Watch him as we pleased, we could do nothing to solve it; and when we asked him to his face, he would only laugh if he were drunk, and if he were sober deny solemnly that he ever tasted anything but water.

He was not only useless as an officer and a bad influence amongst the men, but it was plain that at this rate he must soon kill himself outright, so nobody was much surprised, nor very sorry, when one dark night, with a head sea, he disappeared entirely and was seen no more.

»Overboard!« said the captain. »Well, gentlemen, that saves the trouble of putting him in irons.«

But there we were, without a mate; and it was necessary, of course, to advance one of the men. The boatswain, Job Anderson, was the likeliest man aboard, and though he kept his old title, he served in a way as mate. Mr Trelawney had followed the sea, and his knowledge made him very useful, for he often took a watch himself in easy weather. And the coxswain, Israel Hands, was a careful, wily, old, experienced seaman who could be trusted at a pinch with almost anything.

He was a great confidant of Long John Silver, and so the mention of his name leads me on to speak of our ship's cook, Barbecue, as the men called him.

Aboard ship he carried his crutch by a lanyard round his neck, to have both hands as free as possible. It was something to see him wedge the foot of the crutch against a bulkhead, and propped against it, yielding to every movement of the ship, get on with his cooking like someone safe ashore. Still more strange was it to see him in the heaviest of weather cross the deck. He had a line or two rigged up to help him across the widest spaces—Long John's earrings, they were called; and he would hand himself from one place to another, now using the crutch, now trailing it alongside by the lanyard, as quickly as another man could walk. Yet some of the men who had sailed with him before expressed their pity to see him so reduced.

»He's no common man, Barbecue,« said the coxswain to me. »He had good schooling in his young days and can speak like a book when so minded; and brave—a lion's nothing alongside of Long John! I seen him grapple four and knock their heads together—him unarmed.«

All the crew respected and even obeyed him. He had a way of talking to each and doing everybody some particular service. To me he was unweariedly kind, and always glad to see me in the galley, which he kept as clean as a new pin, the dishes hanging up burnished and his parrot in a cage in one corner.

»Come away, Hawkins,« he would say; »come and have a yarn with John. Nobody more welcome than yourself, my son. Sit you down and hear the news. Here's Cap'n Flint—I calls my parrot Cap'n Flint, after the famous buccaneer—here's Cap'n Flint predicting success to our v'yage. Wasn't you, cap'n?«

And the parrot would say, with great rapidity, »Pieces of eight! Pieces of eight! Pieces of eight!« till you wondered that it was not out of breath, or till John threw his handkerchief over the cage.

»Now, that bird,« he would say, »is, maybe, two hundred years old, Hawkins—they live forever mostly; and if anybody's seen more wickedness, it must be the devil himself. She's sailed with England, the great Cap'n England, the pirate. She's been at Madagascar, and at Malabar, and Surinam, and Providence, and Portobello. She was at the fishing up of the wrecked plate ships. It's there she learned »Pieces of eight,« and little wonder; three hundred and fifty thousand of `em, Hawkins! She was at the boarding of the viceroy of the Indies out of Goa, she was; and to look at her you would think she was a babby. But you smelt powder—didn't you, cap'n?«

»Stand by to go about,« the parrot would scream.

»Ah, she's a handsome craft, she is,« the cook would say, and give her sugar from his pocket, and then the bird would peck at the bars and swear straight on, passing belief for wickedness. »There,« John would add, »you can't touch pitch and not be mucked, lad. Here's this poor old innocent bird o' mine swearing blue fire, and none the wiser, you may lay to that. She would swear the same, in a manner of speaking, before chaplain.« And John would touch his forelock with a solemn way he had that made me think he was the best of men.

In the meantime, the squire and Captain Smollett were still on pretty distant terms with one another. The squire made no bones about the matter; he despised the captain. The captain, on his part, never spoke but when he was spoken to, and then sharp and short and dry, and not a word wasted. He owned, when driven into a corner, that he seemed to have been wrong about the crew, that some of them were as brisk as he wanted to see and all had behaved fairly well. As for the ship, he had taken a downright fancy to her. »She'll lie a point nearer the wind than a man has a right to expect of his own married wife, sir. But,« he would add, »all I say is, we're not home again, and I don't like the cruise.«

The squire, at this, would turn away and march up and down the deck, chin in air.

»A trifle more of that man,« he would say, »and I shall explode.«

We had some heavy weather, which only proved the qualities of the \textit{Hispaniola}. Every man on board seemed well content, and they must have been hard to please if they had been otherwise, for it is my belief there was never a ship's company so spoiled since Noah put to sea. Double grog was going on the least excuse; there was duff on odd days, as, for instance, if the squire heard it was any man's birthday, and always a barrel of apples standing broached in the waist for anyone to help himself that had a fancy.

 \begin{figure}[ph]
\centering
\includegraphics[width=\linewidth]{chap10manathelmimproved}
\caption{The man at the helm was watching the luff of the sail}
\end{figure}

»Never knew good come of it yet,« the captain said to Dr Livesey. »Spoil forecastle hands, make devils. That's my belief.«

But good did come of the apple barrel, as you shall hear, for if it had not been for that, we should have had no note of warning and might all have perished by the hand of treachery.

This was how it came about.

We had run up the trades to get the wind of the island we were after—I am not allowed to be more plain—and now we were running down for it with a bright lookout day and night. It was about the last day of our outward voyage by the largest computation; some time that night, or at latest before noon of the morrow, we should sight the Treasure Island. We were heading S.S.W. and had a steady breeze abeam and a quiet sea. The \textit{Hispaniola} rolled steadily, dipping her bowsprit now and then with a whiff of spray. All was drawing alow and aloft; everyone was in the bravest spirits because we were now so near an end of the first part of our adventure.

Now, just after sundown, when all my work was over and I was on my way to my berth, it occurred to me that I should like an apple. I ran on deck. The watch was all forward looking out for the island. The man at the helm was watching the luff of the sail and whistling away gently to himself, and that was the only sound excepting the swish of the sea against the bows and around the sides of the ship.




In I got bodily into the apple barrel, and found there was scarce an apple left; but sitting down there in the dark, what with the sound of the waters and the rocking movement of the ship, I had either fallen asleep or was on the point of doing so when a heavy man sat down with rather a clash close by. The barrel shook as he leaned his shoulders against it, and I was just about to jump up when the man began to speak. It was Silver's voice, and before I had heard a dozen words, I would not have shown myself for all the world, but lay there, trembling and listening, in the extreme of fear and curiosity, for from these dozen words I understood that the lives of all the honest men aboard depended upon me alone.